% Generated draft: requires mathematical and source review.
\documentclass[11pt]{article}
\usepackage[margin=25mm]{geometry}
\usepackage{fontspec}
\setmainfont{Latin Modern Roman}
\usepackage{amsmath,amssymb,mathtools,unicode-math}
\setmathfont{Latin Modern Math}
\usepackage{enumitem,needspace,xurl,hyperref}
\setlength{\parindent}{0pt}
\setlength{\parskip}{5pt}
\setlength{\emergencystretch}{3em}
\newcommand{\N}{\mathbb{N}}
\newcommand{\Z}{\mathbb{Z}}
\newcommand{\Q}{\mathbb{Q}}
\newcommand{\R}{\mathbb{R}}
\newcommand{\C}{\mathbb{C}}
\newcommand{\F}{\mathbb{F}}
\newcommand{\A}{\mathbb{A}}
\newcommand{\PP}{\mathbb{P}}
\newcommand{\E}{\mathbb{E}}
\newcommand{\eps}{\varepsilon}
\newcommand{\dd}{\,\mathrm{d}}
\begin{document}
\section*{UnsolvedMath ID: 30000726}
\textbf{Source problem code:} catalogue entry
\subsection{Definitions and mathematical statement}
Use finite simple graphs. A graph is \(r\)-connected if it has more than \(r\) vertices and deleting fewer than \(r\) vertices leaves it connected. The intended three prescribed vertices must be pairwise distinct: if \(x=s\) or \(x=t\), the demanded vertex disjointness would be impossible. This distinctness is implicit in ``three vertices'' but not expanded in the extracted formula.

The source asks whether there is a finite-valued function \(f_3:\Z_{\geq1}\to\Z_{\geq1}\) such that for every \(k\geq1\), every \(f_3(k)\)-connected graph \(G\), and every pairwise distinct \(s,t,x\in V(G)\), there exist a simple \(s\)-to-\(t\) path \(P\) and a \(k\)-connected subgraph \(H\leq G\) satisfying
\[
x\in V(H),\qquad V(H)\cap V(P)=\varnothing.
\]
Equivalently, \(G-V(P)\) must contain a \(k\)-connected subgraph through \(x\). The deleted vertices include the endpoints of \(P\). The whole remainder \(G-V(P)\) is not required by the displayed statement to be \(k\)-connected. No path length, size of \(H\), or growth-rate bound on \(f_3\) is specified. The connectivity threshold depends only on \(k\), not on \(G\) or the chosen vertices.

The supplied literature assessment checked 2026-08-21 relates this to Lov\'asz's removable-path question, distinguishes a solved edge-removal weakening from a remaining vertex-deletion problem, and reports the latter open for general \(k\). It does not reproduce a theorem identifying the present subgraph-through-\(x\) statement with the stronger condition that the entire remainder be \(k\)-connected. Nor would edge deletion alone establish the required vertex disjointness. The precise relationship and special-case hypotheses must be checked rather than strengthening the extracted conclusion.
\subsection{Short English statement}
Is sufficiently high connectivity, depending only on \(k\), enough to find a path between two prescribed vertices disjoint from a \(k\)-connected subgraph through a third? The vertices must be distinct, and the source asks for a subgraph of the remainder rather than connectivity of the whole remainder.
\subsection{Sources}
\begin{itemize}
\item Local UnsolvedMath registry, permanent ID 30000726, catalogue entry, \emph{Lov\'asz Removable-Path Conjecture for Three Prescribed Vertices}; the original wording labels it Lov\'asz-related.
\item \emph{Graph Theory} (2008), Oberwolfach Reports, pp. 887--944, DOI \url{https://doi.org/10.4171/owr/2007/16}.
\item Supplied reference ``weaker Lovasz path-removal result'': \url{https://doi.org/10.1016/j.jctb.2007.11.003}.
\item Supplied literature triage checked 2026-08-21; no exact reduction or equivalence with the original removable-path statement is reproduced.
\end{itemize}
\end{document}
